\section{Resultados medidos}

Todas las cifras son de la \dSoc{} en una \dBoard{} con el firmware del commit \texttt{\dFirmwareCommit}, y
salen de los registros guardados como evidencia. No son del ESP32-S3 ni del receptor u-blox de la PCB final.

\subsection{Tiempo de \texttt{geofence\_check()} (PoC-01)}

Media por llamada de un lote de 40, en tres arranques consecutivos. Los tres arranques coinciden
con una diferencia máxima de 0.02~\% y las 40 clasificaciones de cada polígono fueron correctas.

\begin{table}[H]
\centering
\small
\begin{tabularx}{\textwidth}{|r r r X|}
\hline
\tablehead{\textbf{Vértices} & \textbf{Media por llamada} & \textbf{Margen frente a 50~ms} & \textbf{Estado}} \\ \hline
4 & \dBenchUsFour~$\mu$s & \dBenchMarginFour$\times$ & Cumple \\ \hline
16 & \dBenchMsSixteen~ms & \dBenchMarginSixteen$\times$ & Cumple \\ \hline
64 & \dBenchMsSixtyFour~ms & \dBenchMarginSixtyFour$\times$ & Cumple \\ \hline
256 & \dBenchMsTwoFiftySix~ms & \dBenchMarginTwoFiftySix$\times$ & Cumple \\ \hline
1024 & \dBenchMsThousand~ms & \dBenchMarginThousand$\times$ & \textbf{No cumple} \\ \hline
\end{tabularx}
\caption{Tiempo medio de \texttt{geofence\_check()} en la placa.}
\end{table}

El costo es lineal: $T(N) \approx \dBenchFitOffsetUs~\mu\text{s} + \dBenchFitSlopeUs~\mu\text{s}\cdot N$,
con diferencias de a lo sumo 1.5~$\mu$s respecto a las medias. El límite de 50~ms se alcanza cerca de los
\dBenchLimitVertices{} vértices. \textbf{El criterio PoC-01 se cumple para polígonos de hasta unos 800 vértices
y no se cumple con 1024.} Una cerca de una finca, con decenas o unos cientos de vértices, queda dentro. Es
una medida con el contador del kernel y no con un osciloscopio.

\subsection{Primer fix en frío (PoC-07)}

\begin{table}[H]
\centering
\small
\begin{tabularx}{\textwidth}{|c r X|}
\hline
\tablehead{\textbf{Arranque} & \textbf{Tiempo hasta el LED verde} & \textbf{Relación con los \dTtffSpecCold~s de la hoja de datos}} \\ \hline
1 & \dTtffRunA~s & 3.0$\times$ \\ \hline
2 & \dTtffRunB~s & 1.8$\times$ \\ \hline
3 & \dTtffRunC~s & 4.5$\times$ \\ \hline
\end{tabularx}
\caption{Tiempo hasta el primer fix usable tras un arranque en frío, medido con cronómetro.}
\end{table}

La media es de \dTtffMean~s, la mediana de \dTtffMedian~s y la desviación de \dTtffSd~s, con solo tres
muestras. La hoja de datos del \dSoc{} da \dTtffSpecCold~s típico en frío, en cielo abierto, a 25~°C y con
un amplificador de bajo ruido externo con filtro SAW; no se sabe si esta placa cumple esa suposición.
Las pruebas se hicieron en una calle de barrio entre casas, no en cielo abierto, así que la comparación es
orientativa. La incertidumbre de la medida es de unos 2~s. El tiempo medido es hasta un fix con precisión
reportada de 50~m o mejor, no hasta el primer fix a secas.

Con el mismo arranque en frío, esta vez desde la ventana de la prueba estática, \textbf{no hubo ningún
fix en \dWindowColdMinutes{} minutos}. En ese tiempo la lista de satélites del receptor osciló entre
\dWindowColdListedMin{} y \dWindowColdListedMax, con cuatro o más solo el \dWindowColdFourPct~\% del
tiempo, salvo un pico de \dWindowColdPeak{} en los primeros segundos de la búsqueda. Esa lista cuenta
satélites que aparecen en la respuesta del receptor, no satélites con señal utilizable ni usados en el fix.
Una explicación posible, que no se comprobó, es que sin datos de asistencia el receptor tiene que
decodificar la órbita de cada satélite desde su propia señal, y eso exige una señal más fuerte y continua
que la del cielo parcial de esa ventana.

\subsection{Posición GNSS en la prueba estática (PoC-02 y PoC-03)}

La prueba se hizo desde una ventana de primer piso orientada al sur, abierta, con la placa prácticamente
fuera, con el firmware normal y los datos de satélites guardados. La placa estuvo quieta durante los
\dStaticRecordingMinutes{} minutos que duró la grabación. La secuencia fue:

\begin{itemize}
  \item A los \dStaticFirstFixS~s del reinicio llegó el primer fix, con una precisión reportada de
        \dStaticFirstFixAcc~m, que la aplicación descartó por pasar los 50~m. A los \dStaticFenceS~s un
        fix aceptable creó la cerca y la aplicación pasó a \texttt{MONITORING}.
  \item \textbf{Primer tramo:} hasta los \dStaticRunOneEndS~s, \dStaticRunOneFixes{} fixes aceptados a 1~Hz
        (\dStaticRunOneMinutes~min), con un solo hueco de \dStaticMaxGapS~s. La precisión reportada
        convergió de \dStaticAccStart~m a unos \dStaticAccConverged~m en menos de un minuto. Al final se
        degradó de forma gradual, de \dStaticTailAccFrom{} a \dStaticTailAccTo~m, la lista de satélites bajó a
        3 o 4, y a los \dStaticLostFirstS~s la aplicación pasó a \texttt{FIX\_LOST}.
  \item \textbf{Sin fix durante \dStaticGapMinutes{} minutos.} El receptor siguió buscando con la placa quieta.
  \item \textbf{Segundo tramo:} a los \dStaticRunTwoStartS~s el fix volvió, con una precisión reportada
        inicial de \dStaticRunTwoFirstAcc~m, y hubo \dStaticRunTwoFixes{} fixes aceptados hasta los
        \dStaticRunTwoEndS~s (\dStaticRunTwoMinutes~min). A los \dStaticLostSecondS~s se degradó de nuevo y la
        aplicación volvió a \texttt{FIX\_LOST}.
\end{itemize}

Que el fix se pierda y vuelva con la placa quieta indica que lo que cambia es la señal que llega, no la
placa. La causa no se determinó: el firmware de la PoC no registra la intensidad de señal de cada satélite.
Una hipótesis, sin comprobar, es que los satélites, al moverse por el cielo, pasan detrás de los árboles, los
cables o el borde de la ventana.

\textbf{PoC-03.} El intervalo entre fixes dentro de cada tramo fue de 1.01~s: el receptor entrega 1~Hz en modo
continuo, y por tanto \textbf{no supera 1~Hz}, como pide SR-10. Ese requisito apunta al receptor de la PCB.

\textbf{PoC-02.} Se descartaron los primeros \dStaticSkipS~s de fixes de cada tramo, por la convergencia del
receptor, y se calculó la dispersión de las posiciones respecto a su propio promedio. Quedaron
\dStaticBothUsed{} fixes, unos 9.5 minutos de datos, por debajo de los 15 minutos previstos.

\begin{table}[H]
\centering
\small
\begin{tabularx}{\textwidth}{|X r r r|}
\hline
\tablehead{\textbf{Estadístico} & \textbf{Tramo 1} & \textbf{Tramo 2} & \textbf{Ambos tramos}} \\ \hline
Fixes usados & \dStaticRunOneUsed & \dStaticRunTwoUsed & \dStaticBothUsed \\ \hline
CEP50 (m) & \dStaticRunOneCepFifty & \dStaticRunTwoCepFifty & \dStaticBothCepFifty \\ \hline
CEP95 (m) & \dStaticRunOneCepNinetyFive & \dStaticRunTwoCepNinetyFive & \dStaticBothCepNinetyFive \\ \hline
DRMS (m) & \dStaticRunOneDrms & \dStaticRunTwoDrms & \dStaticBothDrms \\ \hline
Precisión reportada, mediana (m) & \dStaticRunOneAcc & \dStaticRunTwoAcc & \dStaticBothAcc \\ \hline
\end{tabularx}
\caption{Dispersión de la posición en la prueba estática, tras descartar los primeros 60~s de cada tramo.}
\end{table}

Con solo los fixes cuya precisión reportada fue de \dStableLimit~m o mejor (\dStableUsed{} fixes de los dos
tramos), el CEP50 es de \dStableCepFifty~m, el CEP95 de \dStableCepNinetyFive~m y el DRMS de \dStableDrms~m;
filtrar por la precisión reportada no mejora la dispersión.

Frente al criterio de \dReqCep~m CEP del SRD, el CEP50 medido es \dStaticRatioTwo{} veces mayor en el mejor tramo,
\dStaticRatioOne{} veces en el peor y \dStaticRatioBoth{} veces con los dos juntos. Frente a la hoja de datos del
\dSoc{}, que da CEP50 de \dNordicCepFifty~m en seguimiento continuo con cielo abierto, tampoco coincide, pero las
condiciones no son las mismas. \textbf{El criterio PoC-02 no se cumple en estas condiciones.} Esto no permite
decir qué daría con cielo abierto.

Tres observaciones sobre los números:

\begin{itemize}
  \item \textbf{La dispersión de los dos tramos juntos es mayor que la de cada uno.} Los promedios de ambos
        tramos están separados por \dStaticRunSeparation~m. Es decir, el error tiene una parte lenta, que
        cambia con la constelación visible, y una prueba corta la subestima. Esto sugiere que el valor de un
        tramo aislado, como el de 3.5~m del segundo, no es representativo.
  \item \textbf{La cola es pesada:} el CEP95 es unas 2.7 veces el CEP50 en cada tramo, es
        decir, la mayoría de los fixes están cerca del promedio pero hay desvíos grandes.
  \item \textbf{La precisión que reporta el receptor} (1 sigma, mediana de unos 11 a 13~m) es mayor que la
        dispersión observada dentro de un tramo (DRMS de 5 a 8~m). No se contradicen: la dispersión alrededor
        del promedio no ve un error constante que afecte a todos los fixes del tramo por igual.
\end{itemize}

\subsection{Alerta del LED y modo demo (PoC-04, PoC-05 y PoC-11)}

El registro del modo demo muestra \texttt{SEARCHING} $\to$ \texttt{DEMO}, la creación de la cerca de 60~m
alrededor del punto ficticio, 42 pasos de la ruta y \texttt{DEMO} $\to$ \texttt{SEARCHING} al pulsar de
nuevo. Algunos pasos, tal como salieron en el registro:

\begin{table}[H]
\centering
\small
\begin{tabularx}{\textwidth}{|r c c X|}
\hline
\tablehead{\textbf{Distancia al borde} & \textbf{Dentro} & \textbf{Color (R, G, B)} & \textbf{LED}} \\ \hline
30.00~m & sí & 0, 255, 0 & Verde fijo \\ \hline
24.00~m & sí & 0, 255, 0 & Verde fijo \\ \hline
18.00~m & sí & 26, 229, 0 & Empieza a mezclarse hacia el rojo \\ \hline
12.00~m & sí & 102, 153, 0 & Mezcla \\ \hline
6.00~m & sí & 179, 76, 0 & Mucho más rojo \\ \hline
0.00~m & no & 255, 0, 0 & Rojo parpadeando (en el borde exacto la regla semiabierta da fuera) \\ \hline
6.00~m & no & 255, 0, 0 & Rojo parpadeando, fuera \\ \hline
4.00~m & sí & 204, 51, 0 & De regreso, dentro \\ \hline
10.00~m & sí & 128, 127, 0 & Mezcla \\ \hline
22.00~m & sí & 0, 255, 0 & Verde fijo \\ \hline
\end{tabularx}
\caption{Pasos del modo demo (\texttt{evidence/logs/demo\_mode\_session.log}).}
\end{table}

La secuencia coincide con el diseño: verde fijo a 20~m o más del borde, mezcla hacia el rojo al acercarse,
y rojo parpadeante fuera, hasta 10~m más allá del borde, y de vuelta. La persona que hizo la prueba pulsó el
botón, vio la placa entrar al modo demo y, al pulsar de nuevo, vio el LED volver al azul parpadeando. El
modo demo funcionó sin señal GNSS. \textbf{PoC-04, PoC-05 y PoC-11 se cumplen} en la PoC.

\subsection{Comportamiento a prueba de fallos (PoC-06)}

\begin{itemize}
  \item En la prueba estática, el registro muestra \texttt{MONITORING} $\to$ \texttt{FIX\_LOST} a los
        \dStaticLostFirstS~s y otra vez a los \dStaticLostSecondS~s, en cada caso unos diez segundos después del
        último fix aceptado, y \texttt{FIX\_LOST} $\to$ \texttt{MONITORING} cuando el fix volvió.
  \item La persona que hizo las pruebas observó el LED en magenta parpadeando al perderse el fix, distinto del
        rojo parpadeante de estar fuera de la cerca.
  \item Los tests del PC comprueban que una distancia no finita o negativa, un \texttt{safe\_m} no positivo y
        cualquier estado desconocido dan el resultado a prueba de fallos, y que \texttt{FAULT} solo se
        abandona con el modo demo.
\end{itemize}

\subsection{Caminata de campo: precisión al aire libre y pérdida de señal (PoC-02 y PoC-06)}

El 20 de noviembre se hicieron dos capturas de una caminata con la placa conectada por USB a un teléfono. La
aplicación del teléfono no guardó el archivo del registro, así que las cifras salen de una \emph{transcripción}
línea por línea del texto copiado, no de una captura cruda (\texttt{field\_\allowbreak walk\_\allowbreak open\_\allowbreak sky.csv} y
\texttt{field\_\allowbreak walk\_\allowbreak signal\_\allowbreak loss.csv}, sin coordenadas). Los segundos de la primera están completos y
consecutivos.

\textbf{Al aire libre.} Durante \dWalkSeconds~s (\dWalkFixes{} fixes, uno por segundo) la persona caminó dentro de la
cerca, de \dWalkDistMin{} a \dWalkDistMax~m del borde, con cielo despejado. La precisión que reportó el receptor
fue de \dWalkAccMin{} a \dWalkAccMax~m, con mediana de \dWalkAccMedian~m: \dWalkWithin{} de los \dWalkFixes{} fixes
(\dWalkWithinPct~\%) reportaron \dReqCep~m o menos y los otros \dWalkAbove{} reportaron \dWalkAccMax~m. El color
siguió la distancia al borde y la comprobación tardó \dCrossCheckInUs~\textmu s. Frente a los 11 a 13~m
reportados en la prueba estática con cielo parcial, la diferencia es grande, y muestra cuánto depende la
precisión de la vista al cielo.

Estas cifras \textbf{no demuestran PoC-02}: la precisión reportada es la estimación del propio receptor (2D, 1
sigma), no el error frente a un punto conocido; son \dWalkSeconds~s en movimiento, no una prueba estática de
15 minutos; y no hubo referencia de coordenadas. Son un indicio favorable de lo que podría dar el receptor con
cielo abierto, que es lo que habría que medir.

\textbf{Cruce del borde y pérdida de señal.} En la otra captura, caminando hacia el borde, el LED pasó de
verde a rojo y la distancia pasó de 1.20~m dentro a 0.31~m fuera, con el parpadeo activado y la comprobación en
\dCrossCheckOutUs~\textmu s. Según la persona que hizo la prueba, al entrar a una casa se perdió la señal. El
registro lo muestra así:

\begin{itemize}
  \item La precisión reportada de los \dCrossFixes{} fixes aceptados creció de \dCrossAccFrom{} a \dCrossAccTo~m, y los
        \dCrossIgnored{} siguientes (hasta \dCrossIgnoredLast~m) se descartaron por superar 50~m.
  \item \texttt{MONITORING} $\to$ \texttt{FIX\_LOST} a los \dCrossLostS~s del último fix aceptado, como en la
        prueba estática (PoC-06).
  \item Los satélites en seguimiento bajaron de \dCrossSatsFrom{} a \dCrossSatsTo{} en \dCrossSatsSeconds~s, y
        quedaron en 0 al final de la captura.
  \item \textbf{Una observación sin explicar:} en los últimos \dCrossConstSteps{} fixes aceptados el avance entre
        fixes fue constante, de \dCrossConstLow{} a \dCrossConstHigh~m por segundo, mientras que antes variaba de
        \dCrossVarLow{} a \dCrossVarHigh~m. Una explicación posible es que el receptor siguiera estimando la
        posición con la última velocidad al perder los satélites. No está comprobado: el registro no dice
        cuántos satélites se usaron en cada fix.
\end{itemize}

Si esa explicación fuera cierta, el firmware habría mostrado el rojo del borde con una posición que ya no se
medía, porque acepta fixes con precisión de hasta 50~m. Por eso el diagnóstico de la señal figura entre los
trabajos futuros.

\textbf{Final de una prueba estática de unos 15 minutos.} Ese mismo día la persona se quedó quieta unos 15 minutos y
al final volvió caminando a una casa. El monitor serie del teléfono solo conservó el último minuto, así que
\textbf{los primeros minutos de esa prueba no se registraron} y no se puede calcular un CEP de ella. Lo conservado
(\texttt{field\_\allowbreak static\_\allowbreak tail.csv}, transcrito) son \dTailFixes{} fixes en \dTailSpanS~s, todos fuera de la cerca:

\begin{itemize}
  \item En los primeros \dTailStillFixes{} fixes, con la persona quieta a unos 341~m del borde, la precisión reportada
        fue de \dTailStillAccLow{} a \dTailStillAccHigh~m y la distancia al borde varió de \dTailStillDistLow{} a
        \dTailStillDistHigh~m. Junto con la caminata al aire libre (\dWalkAccMin{} a \dWalkAccMax~m), es otro indicio de
        que con cielo abierto el receptor estima entre 2 y 3~m. Sigue siendo la estimación del receptor, no el
        error frente a un punto conocido.
  \item Después, mientras la persona volvía a la casa, la precisión reportada creció de \dTailStillAccHigh{} a
        \dTailAccTo~m en \dTailDegradeS~s y la posición se desplazó unos \dTailShiftM~m. Los siguientes \dTailIgnored{}
        fixes (de \dTailIgnoredLow{} a \dTailIgnoredHigh~m) se descartaron y \texttt{MONITORING} $\to$ \texttt{FIX\_LOST}
        llegó \dTailLostS~s después del último fix aceptado, la cuarta vez que se observa. Los satélites en
        seguimiento eran \dTailSats{} en ese instante, lo que sugiere (no está comprobado) que el conteo se actualiza con retraso y que la precisión
        reportada avisa antes.
  \item El LED siguió en rojo parpadeante todo el tiempo y la comprobación tardó \dTailCheckUs~\textmu s.
\end{itemize}

Aquí el desplazamiento de la posición se explica porque la persona caminaba, así que esta captura no apoya ni
descarta la hipótesis de extrapolación de la caminata anterior.

\subsection{Huella de memoria (PoC-08)}

El firmware normal ocupa \dFlashBytes~B de flash (\dFlashPercent~\% de los \dFlashRegion~B disponibles para
la aplicación) y \dRamBytes~B de RAM (\dRamPercent~\% de \dRamRegion~B). La compilación con el benchmark ocupa
\dFlashBench~B y la que borra los datos de satélites, \dFlashCold~B; ambas usan la misma RAM. El motor de
cerca ocupa \dEngineBytesCmThree~B compilado para Cortex-M33 y \dEngineBytesEsp~B para ESP32-S3.

\subsection{Corrección del algoritmo y portabilidad (PoC-09 y PoC-10)}

El motor pasa sus \dHostTests{} pruebas (\dHostChecks{} comprobaciones) y \dStressCases{} casos aleatorios con semilla
fija, sin fallos, también con detectores de errores de memoria. El peor error de la distancia del motor
frente a haversine fue de \dStressWorstErrorPct~\%, diez veces menor que el 0.2~\% documentado. \textbf{PoC-09
se cumple.}

Los cinco módulos de C puro y el motor compilan sin cambios, con las advertencias estrictas como errores, para
el PC, el Cortex-M33 y el ESP32-S3. \textbf{PoC-10 se cumple como compilación}; no se ejecutó nada en el
ESP32-S3.
